\section*{\Huge Résumé}

\vspace{0.7cm}

\begin{spacing}{1.25}
Ce rapport présente le travail réalisé dans le cadre de mon Projet de Fin d’Études, portant sur la conception, le développement et le déploiement de FinderrLink, une plateforme web de mise en relation entre freelances, entreprises et agences. Le projet vise à faciliter la recherche de missions, la gestion des profils professionnels, le suivi des candidatures et la communication entre les différents acteurs de la plateforme.

FinderrLink repose sur une architecture full-stack moderne composée d’un frontend dynamique, d’un backend structuré et d’une base de données relationnelle. La plateforme intègre plusieurs modules fonctionnels, notamment la gestion des utilisateurs, des profils freelances, des missions, des candidatures, de la messagerie, des notifications, des tableaux de bord et d’un assistant intelligent destiné à améliorer l’expérience utilisateur.

Un élément important du projet est la mise en place d’un système RBAC, ou Role-Based Access Control, permettant de gérer les droits d’accès selon le rôle de chaque utilisateur. Les freelances, les entreprises, les agences et les administrateurs disposent ainsi d’interfaces, de permissions et de fonctionnalités adaptées à leurs besoins. Ce système permet d’assurer une meilleure sécurité, une séparation claire des responsabilités et une gestion cohérente des accès.

Le projet couvre également la phase de déploiement dans un environnement serveur, incluant la configuration des services applicatifs, la gestion de la base de données, la mise en place du reverse proxy, la sécurisation des accès et la préparation de l’application pour une utilisation en production.
\end{spacing}

\vspace{0.4cm}
\noindent \textbf{Mots clés :} FinderrLink, Développement full-stack, Déploiement web, Plateforme web, Freelance, Entreprise, Agence, RBAC, Missions, Candidatures, Messagerie, Notifications, Intelligence artificielle, API, Base de données.
